\documentclass[10pt,a4paper]{amsart}
\usepackage[margin=19mm]{geometry}
\usepackage{amsmath,amssymb,mathtools}
\usepackage[scr=rsfs,cal=euler]{mathalfa}
\usepackage{xfrac}
\usepackage[unicode,hidelinks,bookmarks=false]{hyperref}
\newcommand{\ie}{i.e.\ }
\newcommand{\cf}{cf.\ }
\newcommand{\resp}{respectively\ }
\newcommand{\Subsection}{\subsection}
\providecommand{\nobreakdash}{\nobreak\hbox{-}}
\setlength{\emergencystretch}{2em}
\multlinegap0pt
\usepackage{mathtools}
\usepackage{amssymb}
\usepackage{cancel}
\newtheorem{proposition}{Proposition}
\title{An archimedean approach to singular moduli on Shimura curves}
\author{Mateo Crabit Nicolau}
\date{Source: arXiv:2603.08569v1 (CC BY 4.0)}
\begin{document}
\maketitle

\noindent\emph{Source and licence.} An archimedean approach to singular moduli on Shimura curves. Version arXiv:2603.08569v1 (CC BY 4.0), distributed by arXiv under the Creative Commons Attribution 4.0 International licence, http://creativecommons.org/licenses/by/4.0/, as recorded in the arXiv licence metadata for this exact version and shown on the abstract page https://arxiv.org/abs/2603.08569v1. Full licence notice: \texttt{licenses/arxiv-2603.08569-CC-BY-4.0.txt}.\par\smallskip

\noindent\textit{Editorial scope.} Counted unit: the article's proposition Greenf (source0.tex lines 332--340). The article's complete original proof of it is delivered with it (source0.tex lines 341--350, one \texttt{proof} environment of 10 lines). No mathematics is written, repaired or re-derived: the statement and the proof are the article's own lines, in the article's own notation, and no retained line is substituted or re-flowed.  No further source-local context is needed: every cross-reference of every retained line resolves inside the two delivered spans.  Nothing was omitted from any retained span, so the package carries no omission record.  No retained line contains a citation, so the wrapper carries no bibliography and no bibliography entry.  No retained line contains a figure environment, an image-inclusion command or drawing code, so no image is delivered and the package declares no asset dependency.  Editorial formatting only, in addition: an AMS \texttt{amsart} preamble that restates the article's own declarations and macros used by the retained lines, namely the environments \texttt{proposition} on separate counters, together with the standard packages those lines need.  The retained environments print in this document as Proposition 1 in order of appearance, whereas the article prints the same items with the numbering of its own full document.  The complete native source of the article as distributed in the arXiv e-print for this exact version is bundled unchanged as \texttt{sources/196015/source0.tex}.
\begin{proposition}\label{Greenf}
    Let $\tau_i\in \mathbf{H}$ corresponding to a CM-point associated to $K_i$ for $i=1,2$. Then
    \begin{equation}
    \begin{split}
        \log(\mid J_N(\tau_1,w_q(\tau_1),\tau_2,w_p(\tau_2)\mid^2)&=G_N(\tau_1,\tau_2)+G_N(w_q(\tau_1),w_p(\tau_2))
        \\&-G_N(w_q(\tau_1),\tau_2)-G_N(\tau_1,w_p(\tau_2)
        \end{split}
    \end{equation}
\end{proposition}

\begin{proof}
    The function 
    \begin{equation}
        f:z\rightarrow \log\Biggl(\frac{\mid j_{N}(z)-j_{N}(\tau_2)\mid^2}{\mid j_{N}(z)-j_{N}(w_p(\tau_2))\mid^2}\Biggl)
        -(G_N(z,\tau_2)-G_N(z,w_p(\tau_2)))
    \end{equation}
is harmonic and defined over $\mathbf{H}$ even when $z=\tau_2 \textup{ or } w_p(\tau_2) \mod R^1$ because both sides have the same logarithmic pole here. As $R^1\backslash\mathbf{H}$ is compact, this function is constant. In particular, we have
\[f(\tau_1)-f(w_q(\tau_1))=0\] 
which conclude the proof.
\end{proof}

\end{document}
